\chapter{About this manual}

LinkStudio is a block-based simulator of the physical layer (PHY) of communication systems. A
simulation is a directed graph of blocks: sources, modulators, pulse shapers, channels,
synchronizers, equalizers, decoders and meters. The same engine covers radio links,
copper lines and optical fiber. It is driven from Python, from YAML scenario files, from the
command line or from a browser-based editor.

This manual has two goals and is organized accordingly.

\begin{description}[leftmargin=0pt,font=\sffamily\bfseries\color{accent}]
  \item[Part I -- Using and extending LinkStudio.] How the engine works (signals, blocks, graphs,
  Monte-Carlo runs, sweeps, scenarios) and how to drive it from the GUI, the CLI and the
  Python API. It ends with a developer guide for writing new blocks.
  \item[Part II -- Theory and implementation.] One chapter per subsystem. Every chapter first
  develops the theory with the notation used in the code, then shows \emph{where} and
  \emph{how} it is implemented, and finally validates the implementation against a closed-form
  result. Part II can be read as a compact textbook on digital transmission whose equations
  can all be checked in the code.
\end{description}

\section*{How to read the chapters of Part II}
Each theory chapter follows the same pattern:
\begin{enumerate}
  \item \textbf{Model.} The physical or mathematical model, derived with the assumptions stated
  explicitly.
  \item \textbf{Implementation.} A blue \emph{In the code} box names the blocks, files and
  metadata involved. Short excerpts show the lines that realize the key equations, for example
  the split-step solver of the nonlinear Schrödinger equation (\cref{lst:ssfm}). Excerpts are
  extracted automatically from the source, so their file names and line numbers match the code
  base this manual was built from.
  \item \textbf{Validation.} A figure comparing the simulator with theory. All measured curves
  in this manual (markers in the plots) are produced by running the real blocks with
  \file{manual/tools/generate.py}; solid lines are closed-form results.
\end{enumerate}

Green \emph{Key idea} boxes summarize a result; orange \emph{Watch out} boxes point at
common mistakes or at modelling limitations.

\section*{Typographic conventions}
\begin{center}
\begin{tabular}{@{}ll@{}}
\toprule
\blk{awgn} & a block type (as written in scenarios and in the GUI palette)\\
\prm{snr_db} & a block parameter\\
\meta{noise_var} & a metadata key travelling with a signal\\
\file{linkstudio/phy/optical.py} & a file, relative to \file{backend/}\\
\kind{SAMPLES} & a signal kind (\cref{sec:signal})\\
\bottomrule
\end{tabular}
\end{center}

\section*{Notation}
\begin{center}
\begin{tabular}{@{}ll@{\qquad}ll@{}}
\toprule
$x(t)$, $x[n]$ & complex envelope (continuous, sampled) & $f_s$ & sample rate\\
$T$, $R_s=1/T$ & symbol period, symbol rate & $N_{\mathrm{sps}}$ & samples per symbol\\
$E_s$, $E_b$ & energy per symbol, per information bit & $N_0$ & one-sided noise PSD\\
$M$, $k=\log_2 M$ & constellation size, bits per symbol & $R$ & code rate\\
$\CN(\mu,\sigma^2)$ & circular complex Gaussian & $\Q(x)$ & Gaussian tail function\\
$(\cdot)\conj$, $(\cdot)\T$, $(\cdot)\Her$ & conjugate, transpose, Hermitian & $\star$ & convolution\\
$\F\{\cdot\}$ & Fourier transform & $\E[\cdot]$ & expectation\\
\bottomrule
\end{tabular}
\end{center}
Fourier transforms follow the NumPy sign convention used everywhere in the code:
\begin{equation}
  X(\omega)=\int x(t)\,e^{-j\omega t}\,dt, \qquad x(t)=\frac{1}{2\pi}\int X(\omega)\,e^{j\omega t}\,d\omega .
  \label{eq:fourier-convention}
\end{equation}
This convention fixes the signs of the dispersion operator and of the nonlinear phase in the
fiber model (\cref{sec:nlse}). The discrete transforms are \code{np.fft.fft} and
\code{np.fft.ifft}; they are unitary (\code{norm="ortho"}) in OFDM and DMT.

\section*{Rebuilding this manual}
\begin{lstlisting}[style=shell]
python manual/tools/generate.py      # run the simulator, extract excerpts and tables
cd manual && latexmk -pdf linkstudio-manual.tex
\end{lstlisting}
Re-running the generator after a code change updates every figure, excerpt and reference
table, so the manual stays consistent with the code.
